\documentclass[compress,10pt,xcolor=dvipsnames]{beamer}
\usepackage{pgfpages}
\setbeameroption{show notes on second screen}
\setbeamerfont{note page}{size=\footnotesize}

\usepackage[english]{babel}
\usepackage[utf8]{inputenc}
\usepackage{times}
\usepackage[T1]{fontenc}
\usepackage{alltt}

\usepackage{tikz}
\usetikzlibrary{calc,arrows,positioning,decorations.markings,shapes}
\usepackage{amsfonts}
\usepackage{amssymb}
\usepackage{amsmath}
\usepackage{graphicx}
\usepackage{tabularx}
\usepackage{hyperref}
\usepackage{natbib}
\setbeamertemplate{bibliography item}[text]

%% COLORS
\definecolor{awesome}{RGB}{3, 149, 222}
\definecolor{high1}{RGB}{3, 149, 222}
\definecolor{high2}{HTML}{F58137}
\definecolor{high3}{RGB}{120, 81, 169}
\definecolor{lightergrey}{rgb}{0.9, 0.9, 0.9}

%% TIKZ NODES & STYLES
\tikzset{
  boxnode/.style={draw=awesome!80, rectangle, rounded corners=2pt, minimum width=2.2cm, minimum height=1.0cm, align=center, thick},
  arrow/.style={>=stealth, thick, <->}
}

\usefonttheme{professionalfonts}
\usefonttheme{serif}

\mode<presentation>
{
    \usenavigationsymbolstemplate{}
    \setbeamertemplate{headline}
    {
        \begin{beamercolorbox}[sep=1ex,wd=\paperwidth]{whitebox}
            \hfill{\color{awesome}\insertframenumber}\hspace*{0.4em}
        \end{beamercolorbox}
    }

    \setbeamercovered{invisible}
    \setbeamercolor{frametitle}{fg=awesome}
    \setbeamercolor{date in head/foot}{fg=awesome}
    \setbeamercolor{itemize item}{fg=awesome}
    \setbeamercolor{itemize subitem}{fg=awesome}
    \setbeamercolor{itemize subsubitem}{fg=awesome}
    \setbeamercolor{enumerate item}{fg=awesome}

    \setbeamertemplate{itemize item}[square]
    \setbeamertemplate{itemize subitem}[circle]
    \setbeamertemplate{itemize subsubitem}[triangle]
}

\setbeamertemplate{footline}[text line]{%
  \parbox{\linewidth}{\vspace*{-6pt}{\color{lightergrey}\copyright Mohammad Abdulaziz}\hfill\insertshortauthor\hfill{\color{awesome}}}}
\setbeamertemplate{navigation symbols}{}

\newcommand{\chictitle}[2]{
  \begin{frame}[plain,noframenumbering]
    \begin{center}
      \vspace{2.5em}
      {\color{awesome}\Large\textbf{#1}} \\
      \vspace{0.8em}
      \small{#2}
      \vspace{3em}
      \begin{center}
        {\color{awesome}\large Mohammad Abdulaziz}
      \end{center}
    \end{center}
  \end{frame}
}


\setbeamersize{text margin left=8pt,text margin right=8pt}

\title[Program Verification & Loop Invariants]{Program Verification & Loop Invariants}
\author[7CCSACTP]{Mohammad Abdulaziz}
\date{Week 9}

\begin{document}

\chictitle{Program Verification & Loop Invariants}{Computer-Assisted Theorem Proving (7CCSACTP)}

\begin{frame}{1. The Need for Invariants in Program Verification}
\note{
\begin{itemize}
  \item From pure structural induction to iterative state transformations
  \item State spaces, transition systems, and reachable states
  \item Formulating inductive invariants
\end{itemize}
}

\begin{itemize}
  \item \textbf{From pure structural induction to iterative state transformations:} From pure structural induction to iterative state transformations
  \item \textbf{State spaces, transition systems, and reachable states:} State spaces, transition systems, and reachable states
  \item \textbf{Formulating inductive invariants:} Formulating inductive invariants
\end{itemize}
\end{frame}

\begin{frame}{2. Loop Invariants for Graph & Search Algorithms}
\note{
\begin{itemize}
  \item DFS loop invariants: visited sets, frontier, and target states
  \item Preservation of invariants across traversal steps
  \item Inductive proofs over state transitions
\end{itemize}
}

\begin{itemize}
  \item \textbf{DFS loop invariants:}
    \begin{itemize}
      \item visited sets, frontier, and target states
    \end{itemize}
  \item \textbf{Preservation of invariants across traversal steps:} Preservation of invariants across traversal steps
  \item \textbf{Inductive proofs over state transitions:} Inductive proofs over state transitions
\end{itemize}
\end{frame}

\begin{frame}{3. State Invariants & Safety Properties}
\note{
\begin{itemize}
  \item Proving invariant adequacy: invariant implies postcondition
  \item Handling non-trivial invariant generalisations
  \item Bridge to imperative Hoare logic and weakest preconditions
\end{itemize}
}

\begin{itemize}
  \item \textbf{Proving invariant adequacy:}
    \begin{itemize}
      \item invariant implies postcondition
    \end{itemize}
  \item \textbf{Handling non-trivial invariant generalisations:} Handling non-trivial invariant generalisations
  \item \textbf{Bridge to imperative Hoare logic and weakest preconditions:} Bridge to imperative Hoare logic and weakest preconditions
\end{itemize}
\end{frame}


\end{document}
